\subsection{The exact Euler cube on the source shear fans}
\begin{statement}{P3K.66}{Proved}
Retain the actual selected-cusp positive \(\rho\) loop of
P3F.54/P3F.56, its fixed scalar and finite caps, its physical \(\delta\)
action, and the physical vertical-edge fans
\(T_n=T(e_1+ne_2,e_2)\), for all integers \(n\).
In the native fixed-star Cox product the exact integer Jin
characteristic pair and index are
\[
 k_n=6-18n(n-1),\qquad \ell_n=0,\qquad
 A_n=1-3n(n-1).
\]
The native fixed-star free pair is
\((0,-6+18n(n-1))\). In the actual ordinary plumbing product of
P3K.55 the remaining diagonal entries are
\[
 V_{00}=-18n(n-1),\qquad
 \delta_0=-12+18n(n-1).
\]
The proof extends the prescribed boundary line isomorphism through
the actual completed quotient piece and cap, then computes its
oriented cubic integral. It does not select an integer from a
residue modulo \(28\).
\end{statement}
\begin{proof}
\par

\medskip\noindent\textbf{EC.1. The real connection, the actual complex line, and orientation}

\par\nobreak

We first fix the signs of a real circle connection explicitly. Let
\(\mathcal L_\delta\) be the complex line whose nonzero vectors are the
actual physical principal circle bundle. Its local fibre coefficient
is \(V=|V|\exp(2\pi i d)\); the physical \(\delta\) action increases \(d\).
A real principal connection \(A\), normalized by
\(A(\partial_d)=1\), gives the unitary line connection \(2\pi i A\)
in local unit frames. Consequently
\[
 c_1(\mathcal L_\delta)
       =\frac{i}{2\pi}F_{\nabla_\delta}=-dA.
                                                        \tag{EC.1}
\]
For example, in the actual Cox coordinate
\(q_H=(1+jZ)(c_0+j\overline{c_1})\), the physical action is
\(q_H\mapsto q_H e^{-2\pi i t}\). Its real connection is minus the
usual contact connection in that right-complex spinor. Formula
(EC.1) gives the positive meridian class retained in F56 and K42.
It is the same complex line as the source line on a regular source
slice: K6's literal page restriction is
\(x|_{\mathrm{page}}=c_1(M)\). We retain this arrow, rather than change
the associated representation of the circle.

It is convenient to perform the connection calculation on the inverse
line
\[
 \mathcal U=\mathcal L_\delta^{-1},\qquad
 \widehat x=c_1(\mathcal U)=-x=dA.
                                                        \tag{EC.2}
\]
Thus the real forms in SQ.12--15 represent \(\widehat x\) when the
principal coordinate is the increasing physical \(\delta\) coefficient
\(d\). This makes the associated-line sign explicit. The vanishing
Pontryagin pairing of K63 is unaffected by replacing \(x\) by \(-x\);
its statement and relative-coordinate conclusions are unchanged.

We also keep the contracting suspension coefficient
\(\gamma=\gamma_B\). In physical coordinates
\(V=r\exp(2\pi i d)\), the complex threefold has orientation
\(\Omega_S\wedge dr\wedge dd\). Taking the physical \(\delta\) circle
first gives orbit orientation
\(-\Omega_S\wedge d\log r\). The contracting coefficient increases
as \(\log r\) decreases. Hence its orbit orientation is
\(\Omega_S\wedge d\gamma_B\), as calculated in GO.1. The physical
positive-parameter completion therefore uses
\(\Omega_S\wedge d\gamma_B\wedge d\theta\). K42's actual parameter
reversal to the Jin completion reverses this orientation and carries
the same line \(\mathcal L_\delta\).

Write \(r_q\) for an increasing radial coordinate in the q disc and
\(\alpha=\arg(q)/(2\pi)\). The positive source orientation near its
regular boundary is
\[
 dr_q\wedge d\alpha\wedge du\wedge dv .
                                                        \tag{EC.3}
\]
Indeed the source fibre is \(u+\tau v\), with
\(\operatorname{Im}\tau>0\). The two reversals, of the Jin
orientation and of the odd cube \(x^3=(-\widehat x)^3\), cancel.
We can therefore compute the actual Jin integer \(k\) by integrating
the cube of the curvature of \(\mathcal U\) in the auxiliary
orientation
\[
 \Omega_+=dr_q\wedge d\alpha\wedge du\wedge dv
                         \wedge d\gamma_B\wedge d\theta .
                                                        \tag{EC.4}
\]
This is an equality of the specified oriented integrals. No integer
lift of a residue is used to choose its sign.

\par

\medskip\noindent\textbf{EC.2. Full periods and exact boundary line coordinates}

\par\nobreak

Retain SQ.1--2 in full:
\[
 h:(Z,V)\mapsto(\kappa qZ,\alpha_L V),\qquad
 B:(Z,V)\mapsto(gZ,\eta qV),
\]
\[
 Z_n=Z,\quad V_n=Z^{-n}V,\quad h_n=hB^n,\quad B_n=B,
\]
\[
 Q_n=\kappa qg^n,\quad
 \alpha_{L,n}=\alpha_L\eta^n\kappa^{-n}g^{-n^2},
 \qquad\eta_n=\eta g^{-n}.
                                                        \tag{EC.5}
\]
The symbol \(\alpha_L\) here denotes the line multiplier; it is
not the angle \(\alpha\). The \(\rho\) windings of
\((\kappa,\alpha_L,g,\eta)\) are \((-6,0,0,0)\).
The actual periodic logarithms of every zero-winding factor are
retained.

Here is a precise connection choice that retains them. On the regular
boundary, use the actual real period coefficients
\[
 (z,\ell)=u(1,0)+v(\tau,\mu)+d(0,1)+\gamma(G,\beta),
 \qquad z=\frac{\log Z}{2\pi i},\quad
        \ell=\frac{\log V}{2\pi i},
                                                        \tag{EC.6}
\]
where
\(\tau=\log(\kappa q)/(2\pi i)\),
\(\mu=\log\alpha_L/(2\pi i)\),
\(G=\log g/(2\pi i)\) and
\(\beta=\log(\eta q)/(2\pi i)\).
The imaginary parts give exactly the invertible matrix in SQ.4a;
the real parts then determine \(u,d\). In particular
\[
 d=\operatorname{Re}\ell-\operatorname{Re}\mu\,v
                           -\operatorname{Re}\beta\,\gamma.
                                                        \tag{EC.7}
\]
Differentiating (EC.7) retains the differentials and values of both
unit logarithms. Thus using \(d\) as a coefficient coordinate does
not set their phases or their derivatives to zero.

The q and \(\rho\) coordinate changes are respectively
\[
 (u,d)\mapsto(u-v,d-\gamma),\qquad
 (u,d)\mapsto(u+6v,d).
                                                        \tag{EC.8}
\]
The complete zero-winding logarithms are single-valued and introduce
no further integer change. Consequently
\[
 A_0=dd+\alpha\,d\gamma
                                                        \tag{EC.9}
\]
is an actual global real principal connection on this boundary.
In physical \(V\) coordinates it contains all terms obtained by
differentiating (EC.7). It is a permissible connection choice on the
actual bundle, not an omission of its units.

Substitution in (EC.6), including both complex coordinates, gives
exactly
\[
 u_n=u,\quad v_n=v,\quad\gamma_n=\gamma-nv,\quad d_n=d-nu.
                                                        \tag{EC.10}
\]
In the new coordinates the corresponding real connection is
\[
 A_n=dd_n+\alpha\,d\gamma_n+6n\theta\,dv.
                                                        \tag{EC.11}
\]
Its \(\rho\) change is \(d_n\mapsto d_n-6nv\), and its q change is
\(d_n\mapsto d_n-\gamma_n\), so (EC.11) is invariant under the
literal two deck arrows. Equations (EC.9)--(EC.11) give the exact
global difference
\[
 \beta_n=-n\{du+(\alpha-6\theta)dv\}.
                                                        \tag{EC.12}
\]
It is invariant under both changes (EC.8). If one uses the source
coframe \(du+(\alpha-6\theta+\phi_\kappa)dv\), then the additional
one-form \(n\phi_\kappa dv\) is present in (EC.12). It has not been
deleted. Formula (EC.7) explains the same fact for every other unit.

\par

\medskip\noindent\textbf{EC.3. The relative Cech line isomorphism}

\par\nobreak

Let \(\mathcal C\) be the six-dimensional family of completed cusp
quotient pieces for \(T_0\), over the positive \(\theta\) circle. Use
SQ.1--2 to identify each actual \(\mathcal C_n\) with this fixed
smooth piece. Those arrows are periodic, preserve the outer real
period marking, and retain the typed Cartier spinor map and the
physical Hopf cap. Its boundary is denoted \(\mathcal B\).

Choose the graph map \(b:\mathcal C\to\mathbb R/\mathbb Z\) of
SQ.3, with boundary value \(v\). The same proof modifies it by a
supported real function to be constant near the fixed tube and
extends it over the cap. It is periodic in \(\theta\). Define an actual
complex line \(\mathcal T\) on \(\mathcal C\) by cutting \(\theta\)
and gluing its coefficient by
\[
 (\theta,p,w)\longmapsto
       (\theta+1,p,e^{2\pi i b(p)}w).
                                                        \tag{EC.13}
\]
In these cut unit frames its connection is
\(\nabla_{\mathcal T}=d-2\pi i\theta\,db\). The frame change at
\(\theta\) one is the inverse of the displayed coefficient change, so
this connection descends exactly. Its Chern curvature is
\[
                         \Xi=d\theta\wedge db,\qquad \Xi^2=0.
                                                        \tag{EC.14}
\]

The actual normalization constant factor
\(\alpha_{L,n}/\alpha_L\) has \(\rho\) winding \(6n\). In the inverse
line \(\mathcal U\), its Chern class is consequently
\(+6n[d\theta]\smile b\). The other classes of the line difference
are zero: the typed SQ.8 Hopf arrow preserves the meridian value;
the regular \(b\times\gamma\) torus has the integral period section
from SQ.3; and the \(\gamma\) times \(\theta\) value is zero because
\(\eta_n/\eta=g^{-n}\) has its periodic logarithm. The integral
Thom, Mayer--Vietoris and Wang calculation SQ.16 has no torsion
and shows that these are all the evaluations. Thus
\[
 c_1(\mathcal U_n)-c_1(\mathcal U_0)
                   =6n\,c_1(\mathcal T)
             \quad\hbox{on }\mathcal C.               \tag{EC.15}
\]
This is also the direct \(\theta\)-cut/graph-cut clutch square: the
physical \(\delta\) coefficient changes by \(-6n b\), so its dual
coefficient changes by \(+6n b\), exactly (EC.13).
The zero-winding logarithms left in (EC.5) are
\[
 nf_\eta-nf_\kappa-n^2f_g,\qquad -nf_g,
                                                        \tag{EC.16}
\]
together with the actual source-flow pullbacks of their extending
units. They are retained in the bundle whose class is (EC.15).

We next prove the stronger relative assertion with its prescribed
outer line arrow:
\[
 J_n:\mathcal U_n\longrightarrow
                    \mathcal U_0\otimes\mathcal T^{6n}.
                                                        \tag{EC.17}
\]
On \(\mathcal B\), this arrow is specified in the doubly cut real
period frames by (EC.10)--(EC.11) and the clutch (EC.13).
Their deck changes agree exactly, so it is a genuine boundary
bundle isomorphism. The assertion is that this particular boundary
arrow extends, including its integral gauge; an absolute equality
of Chern classes alone would not prove that assertion.

Here are the exact groups and the only possible relative obstruction.
On \(\mathcal B\), the q monodromy has commutator \(u\) with the
source graph loop \(v\); its coefficient is one. The \(\rho\) monodromy
has commutator \(6u\), and all other retained circle generators
commute after abelianization. Hence
\[
 H_1(\mathcal B;\mathbb Z)
       =\mathbb Z\langle\alpha,\theta,v,\gamma\rangle.
                                                        \tag{EC.18}
\]
The \(u\) generator has been killed primitively, so no finite
summand has been suppressed. SQ.16 and its periodic graph and
\(\gamma\) maps give
\[
 H^1(\mathcal C;\mathbb Z)
       =\mathbb Z\langle d\theta,db,d\gamma\rangle,
\quad
 \operatorname{im}H^1(\mathcal C)
       =\mathbb Z\langle d\theta,dv,d\gamma\rangle
                   \subset H^1(\mathcal B).           \tag{EC.19}
\]
The \(\gamma\) map extends through the cap by its real lift near the
fixed sphere, exactly as for \(b\). A section of the \(\theta\)
family outside the fixed tube splits the Wang sequence, so there
is no extra generator or torsion in (EC.19). Therefore the only
boundary winding that does not extend is the q-angle winding.

We test it using an actual disc, rather than a rational cohomology
argument. Choose a constant Tate section \(Z=Z_*\), with
\(|Z_*|=1\), at a smooth normalization point avoiding the
marked zero section, the retained Picard-section divisors and
the two nodal branches. This can be done uniformly: the explicit
Tate points of F54 remain near \(1\) for the chosen small \(\rho\)
circle, whereas a point near \(-1\) stays disjoint from them;
shrink the common q disc once more if necessary. Take the marked \(\gamma\)
value zero. The fixed graph's short arc in SQ.4c is centred at
one half, so this value lies outside it. The section in the
source neighbourhood extends across q zero through a smooth
point of the nodal fibre. Together with this \(\gamma\) value it is
a disc in the orbit piece, disjoint from the fixed tube and its
surgery; it therefore lies in the completed piece as well.
The SQ.2 source horizontal field can be taken to preserve this
section. On the outside collar its real marked-period field is
already zero on \(v=0\); near the node the section is disjoint.
Adding this prescribed zero field to the stated partition of
unity leaves all the node, boundary and cap arrows unchanged.

More generally, at a fixed \(\gamma\) value in this free interval the
suspension-untwisted coordinates give the physical source point
\[
             Z_{\rm phys}=Z_*\exp(\gamma\log g).
                                                        \tag{EC.20}
\]
This is a nonzero extending function on the q disc; its logarithm
uses the actual periodic \(\log g\). The line character arrow
on it is literally multiplication by \(Z_{\rm phys}^{-n}\).
The new untwisting factor
\(\eta_n^\gamma=\eta^\gamma\exp(-n\gamma\log g)\)
cancels the displayed \(g\) factor on passing to the common
untwisted line frame. Its remaining factor is \(Z_*^{-n}\).
Equivalently one can retain both factors in (EC.20); both extend
and are nowhere zero. No fractional power of \(q\) is used in
this ratio: the common \(q^\gamma\) factors cancel before the
ratio is formed. The disc lies in one free suspension interval,
so no B reindex is introduced. The tensor line \(\mathcal T\)
is trivial on this disc at fixed \(\theta\) and \(b=v=0\).
For the chosen value \(\gamma\) zero, the untwisting is literally the
identity: the disc is the ordinary source section with \(V\) in
one fixed nonzero line chart. We can take \(V=1\). In the two
actual first vertex charts this disc and its image have coordinates
\[
 (q/Z_*,Z_*,1)\longmapsto
                 (qZ_*^{\,n-1},Z_*,Z_*^{-n}).
                                                        \tag{EC.20a}
\]
Both charts have physical height equal to the product of their
three coordinates. Their second and third coordinates are nonzero
at q zero, so this is the free locus and no fixed-tube or cap
coordinate is involved. The inverse is given by the opposite
nonzero powers of \(Z_*\). The dual line arrow is multiplication
by \(Z_*^{-n}\). This gives an extending isomorphism
with the unchanged physical \(\delta\) frame. The prescribed boundary
arrow in (EC.17) thus extends over this disc by this explicit line
map. Its q-angle winding is zero.

For completeness, an explicit Cech construction turns these facts
into the full isomorphism. On a good cover of \(\mathcal C\), use
the actual unitary transition functions of
\(\mathcal R_n=\mathcal U_n\otimes\mathcal U_0^{-1}
\otimes\mathcal T^{-6n}\), formed from all the retained maps.
Choose real logarithms on the overlaps. Their Cech coboundary
is an integral two-cocycle. Equation (EC.15) says that cocycle
is the coboundary of an integral one-cochain; subtract that
one-cochain from the chosen logarithms. Our local frame convention
is \(e_j=\exp(2\pi i r_{ij})e_i\). The resulting real
one-cocycle \(r_{ij}\) is solved by
\[
                  a_i=\sum_k\rho_k r_{ki},\qquad
                  a_j-a_i=r_{ij},                    \tag{EC.21}
\]
for a partition of unity subordinate to the cover. Multiplying
the local unit frames by \(\exp(-2\pi i a_i)\) gives a global
unit frame. Since
\(\mathcal R_n=\operatorname{Hom}
(\mathcal U_0\otimes\mathcal T^{6n},\mathcal U_n)\),
this frame is a nowhere-zero map in that direction. Its inverse
is the absolute isomorphism \(J_n\) in (EC.17).
Thus the direction of the required map is fixed by the actual
transitions, including (EC.16).

Compare its boundary value with the prescribed boundary arrow.
Their ratio is a circle-valued function. By the section-disc
calculation its \(\alpha\) coefficient is zero. By (EC.18)--(EC.19)
its other three winding coefficients are integral multiples
of the restrictions of \(\theta\), \(b\), and \(\gamma\). Multiply the
absolute isomorphism by the corresponding extending functions
\(\exp(2\pi i(a\theta+cb+d\gamma))\) to remove these coefficients.
The remaining boundary ratio has a single real logarithm.
Extend that logarithm to the collar, multiply it by a cutoff,
and set it to zero in the remaining interior. Multiplication
by its exponential makes the boundary map exactly the
prescribed one. This proves (EC.17) as a relative bundle map.
It retains every possible boundary integer rather than declaring
the zero-winding unit terms harmless without an extending arrow.

The cap is included in this construction. Its actual line is
the one carried by the SQ.8 Cartier/Hopf map. The supported
change \(b-b'\) is implemented on the \(\theta\) cut by its explicit
exponential in SQ.3. A residual circle gauge on the inner cap
boundary has only its \(\theta\) winding, and this extends over
\(D^3\times S^2\times S^1_\theta\). Thus no additional cap
line class or unspecified framing is inserted in (EC.17).

\par

\medskip\noindent\textbf{EC.4. The affine curvature family and the oriented collar integral}

\par\nobreak

Choose a connection on \(\mathcal U_0\) whose boundary value is
the one just fixed, and write its real Chern curvature as \(F_0\).
On the fixed completed piece use the relative isomorphism
(EC.17) and the tensor connection of (EC.13). Their curvature is
exactly
\[
                         F_n=F_0+6n\Xi .
                                                        \tag{EC.22}
\]
All connections here are on the actual lines. They need not be
Chern connections of the original complex structure; Chern--Weil
evaluation is independent of that choice. The relative map
(EC.17) ensures their specified boundary values agree with
the actual outer gluing. In particular no omitted Mayer--Vietoris
gauge can add a further quadratic or other term.

Keep the original connection on the common exterior, where the
fan change is absent. In piece coordinates its pullback on
\(\mathcal B\) is displayed below. From this point onward, after
pullback by the actual piece map, we write \(\gamma=\gamma_n\)
for the common piece coefficient. The exterior coefficient is
\(\gamma_{\rm ext}=\gamma+nv\), and its \(\delta\) coefficient is
\(d_{\rm ext}=d_n+nu\). Thus the \(\gamma\) in the curvature
formula is the new piece coefficient, not the exterior coefficient.
With this dictionary the literal pullback is
\[
 A_{\mathrm{ext},n}
   =dd_n+\alpha\,d\gamma_n+n(du+\alpha\,dv),
\quad
 F_{\mathrm{ext},n}
   =d\alpha\wedge d\gamma+n\,d\alpha\wedge dv.
                                                        \tag{EC.23}
\]
The difference between the piece and exterior connections is
the global form (EC.12). It has derivative
\[
 d\beta_n=-n\,d\alpha\wedge dv+6n\,d\theta\wedge dv .
                                                        \tag{EC.24}
\]
Directly,
\[
 F_{\mathrm{ext},n}^2=0,\quad (d\beta_n)^2=0,
\quad F_{\mathrm{ext},n}+d\beta_n
            =d\alpha\wedge d\gamma+6n\,d\theta\wedge dv .
                                                        \tag{EC.25}
\]

Orient \(\mathcal B\) as the outward boundary of the piece in
(EC.4):
\[
 \Omega_{\mathcal B}
       =d\alpha\wedge du\wedge dv\wedge d\gamma\wedge d\theta .
                                                        \tag{EC.26}
\]
Each marked angular coordinate has period one. Its integral
over the actual boundary is one: the q/\(\rho\) monodromies shear
\(u\) by integral multiples of \(v\), and hence preserve this
volume form and its unit fundamental cell. The full unit
periods have already entered the coordinate construction
(EC.6); they do not alter this cell or its primitive lattice.

Let \(t\) increase across a gluing collar from the piece toward
the exterior. Its orientation is \(dt\wedge\Omega_{\mathcal B}\).
Interpolate by
\[
 A(t)=A_{\mathrm{ext},n}+\chi(t)\beta_n,\qquad
                  \chi(0)=1,\quad\chi(1)=0.
                                                        \tag{EC.27}
\]
It is constant near each endpoint. The curvature and its cube
are
\[
 F=F_{\mathrm{ext},n}+\chi d\beta_n
                           +\chi' dt\wedge\beta_n,
\]
\[
 F^3=3\chi'\,dt\wedge
          \beta_n\wedge(F_{\mathrm{ext},n}+\chi d\beta_n)^2 .
                                                        \tag{EC.28}
\]
The term without \(dt\) vanishes on the five-dimensional boundary.
The exact exterior-algebra identity is
\[
 3\int_0^1
   \beta_n\wedge(F_{\mathrm{ext},n}+s\,d\beta_n)^2\,ds
                    =18n^2\,\Omega_{\mathcal B}.        \tag{EC.29}
\]
The actual interpolation in (EC.27) runs from \(s=1\) to \(s=0\).
Its contribution is therefore \(-18n^2\), with the sign fixed
by (EC.26)--(EC.28).

The interior contribution is linear because \(\Xi^2=0\):
\[
 F_n^3-F_0^3=18n\,F_0^2\wedge\Xi .
                                                        \tag{EC.30}
\]
The common exterior contributes no change. Its boundary map
has already been accounted for by (EC.23)--(EC.29).
There is consequently one real constant \(L\), independent
of \(n\), such that the actual integer Euler cubes satisfy
\[
                         k_n=k_0+Ln-18n^2 .
                                                        \tag{EC.31}
\]
This conclusion follows from the global relative line map
and the actual oriented gluing. It is not a polynomial assigned
from the local fifth form alone.

\par

\medskip\noindent\textbf{EC.5. Evaluation of the linear term and exact consequences}

\par\nobreak

K42 proves \(k_0=6\). The fan \(T_1\) is literally the reflected
fan of K46: its first triangle is
\((0,e_1+e_2,e_2)\), and translating its second triangle by
\(-e_2\) gives \((0,e_1,e_1+e_2)\). The translation is the
actual B period, including its full unit map, as retained in
K56. K46 constructs the oriented cap and line isomorphism for
that reflected family and proves \(k_1=6\). Thus (EC.31) at
one gives \(L=18\), without any use of a residue.

The exact integer result is
\[
              k_n=6-18n(n-1),\qquad \ell_n=0,\qquad
              A_n=1-3n(n-1).
                                                        \tag{EC.32}
\]
Here \(\ell_n=0\) is the already proved K63 pairing, and
\(24A_n=4k_n-\ell_n\) is the fixed Jin normalization. The
K63 native and ordinary identities now give exact integers:
\[
 (m_0^{\rm native},n_0^{\rm native})
                  =(0,-6+18n(n-1)),
\]
\[
                 V_{00}=-18n(n-1),\qquad
                 \delta_0=-12+18n(n-1).
                                                        \tag{EC.33}
\]
Their reductions agree with K56 and K63. That agreement is a
check on the completed calculation, not its proof of an integer lift.

The exterior-form diagnostic remains preserved unchanged.
The new certificate, \path{code/part3/source_shear_euler_cube.py},
passes 41 exact checks of the actual
period-coordinate arrow, all integral H1 quotients, the line
connection deck invariance, the cubic collar identity with its
endpoint direction, the reflected triangle relabelling, and the
resulting exact polynomial. The geometric Cech, source-disc,
normal and cap arguments above are separate from that algebraic
certificate.
\end{proof}
